\documentclass[11pt]{article}
\usepackage{graphicx} % Required for inserting images

\title{CP\_II - Blatt 2}
\author{Julia Els, 454291 - Felix Krueckel, 454343}
\date{16. April 2025}

% Packages
\usepackage{amsmath}
\usepackage{amssymb}

% Document
\begin{document}

\maketitle

\section*{Aufgabe 1: k-fache Faltung}
\subsection*{a) k $\in \{2\}$}

\subsubsection*{Fall k=2}
Für $k=2$ ist $w = x_1 + x_2$. Wir bezeichnen die Dichte von $w$ als $f_{2}(w)$.
\begin{equation}
    f_{2}(w) = \int_{-\infty}^{\infty} f(x) f(w-x) dx
\end{equation}
Der Integrand $f(x)f(w-x)$ ist nur dann von Null verschieden, wenn beide Terme von Null verschieden sind. Dies bedeutet:
\begin{enumerate}
    \item $0 \le x \le 1$
    \item $0 \le w-x \le 1 \implies w-1 \le x \le w$
\end{enumerate}

\paragraph{Fall 2.1: $0 \le w \le 1$}
In diesem Fall ist $w-1 \le 0$. Die Bedingungen für $x$ werden zu:
$\max(0, w-1) = 0$
$\min(1, w) = w$
Somit ist das Integrationsintervall $[0, w]$.
\begin{equation}
    f_{2}(w) = \int_0^w 1 \cdot 1 dx = \int_0^w dx = [x]_0^w = w
\end{equation}

\paragraph{Fall 2.2: $1 < w \le 2$}
In diesem Fall ist $w-1 > 0$ und $w > 1$. Die Bedingungen für $x$ werden zu:
$\max(0, w-1) = w-1$
$\min(1, w) = 1$
Somit ist das Integrationsintervall $[w-1, 1]$.
\begin{equation}
    f_{2}(w) = \int_{w-1}^1 1 \cdot 1 dx = \int_{w-1}^1 dx = [x]_{w-1}^1 = 1 - (w-1) = 2-w
\end{equation}

\paragraph{Fall 2.3: $w < 0$ oder $w > 2$}
In diesen Fällen ist das Integral Null, da die Schnittmenge der Träger leer ist. \\


Zusammenfassend gilt für $k=2$:
\begin{equation}
    f_{2}(w) = \begin{cases} w & 0 \le w \le 1 \\ 2-w & 1 < w \le 2 \\ 0 & \text{sonst} \end{cases}
\end{equation}

\end{document}